\section{Introduction}
\label{sec_introduction}


\subsection{State of the art}\label{sec_state_of_the_art}

A number of existing studies rely on the hypothesis that the \acl{bflu} surrounding the \acl{mem} is inertialess.
For instance, Arroyo et al. study the interaction between a \acl{mem} and bulk fluid in the Stokes approximation, for simple membrane geometries that can be described by a single parameter \cite{arroyoRelaxationDynamicsFluid2009i}.
% 
Boedec and coauthors combined  \ac{fe} and \ac{be} methods describe an \acl{mem} interacting with  an inertialess \acl{bflu} \cite{boedecIsogeometricFEMBEMSimulations2017}. In there, the membrane surface rheology, such as viscous effects and  flows induced by surface-tension gradients, is neglected.
% 
A \acl{mem} coupled to a \acl{bflu} has been recently studied by Zhou and co-authors  with \ac{be} and \ac{fe} methods, under the assumption that the \acl{bflu} is inertialess and that \acl{mem}  is  inextensible \cite{zhouModelingSimulationOpen2026}.
% 
Barakat and Shaqfeh  modelled vesicles in a circular tube by means of  \ac{be} methods and lubrication theory, by neglecting \acl{mem} rheology and describing \acl{bflu}  as inertialess \cite{barakatStokesFlowVesicles2018}.
% 
Finally, Rodrigues et al., studied the interaction between an inertialess, closed \aclp{mem} and \aclp{bflu}  with \ac{fe} methods, by adopting a minimal description for the fluid, which is solely modelled as a volume constraint \cite{rodriguesSemiimplicitFiniteElement2015}.



Other studies go beyond the inertialess approximation, and rely on simplifying assumptions for the \acl{mem}.
% 
For instance, Bonito et al. propose a \ac{fe} model for a \acl{mem} interacting  an \acl{bflu},  the latter being in the inertial regime  \cite{bonitoDynamicsBiomembranesEffect2011,bonitoParametricFEMGeometric2010}. In this study \acl{mem} is treated as a purely \acl{ela}, characterized  solely by bending rigidity, area and volume constraints---its fluid features, such as tangential and normal flows, are neglected.
% 
Laadhari extended the analyses above to the interaction   between a \acl{mem} and a non-Newtonian \acl{bflu}, in which membrane rheology is still neglected  \cite{laadhariFiniteElementMethodSimulation2023}.
% 
The interaction between a \acl{mem} and a \acl{bflu}---both in the inertial regime---has been studied by Barrett and co-authors, \cite{barrettFiniteElementApproximation2016}, under the assumption that the \ac{mem} is perfectly inextensible. Such analysis is based on  a fluidic approach, in which the \acl{mem} does not obey its own, intrinsic rheology---see below---but it is advected by the surrounding \acl{bflu}. As a result, the membrane tangential and normal velocities are not determined as the result of the surfacial \ac{ns} equations, but they match the respective components of the \acl{bflu} velocity.

\vspace{0.5cm}\\

In this study, we propose a novel \ac{fe} study of the interaction between a \acf{bflu} and a \acf{mem}, based on the \ac{ale} method \cite{malvernIntroductionMechanicsContinuous1969,howellp.AppliedSolidMechanics2008}. Our analysis  includes a complete description of the physical features of both \ac{mem} and \ac{bflu}, it may describe the inertial regime of both, and it accounts for a detailed description of \ac{mem} elasticity, rheology, and extensibility. Unlike fluidic approaches where the \ac{mem} dynamics is determined by the \ac{bflu}, here both the flows and shape of \ac{mem} are determined by intrinsic physical evolution equations. In addition, our study can account for both closed and open surfaces, by leveraging the flexibility of its \ac{fe} formulation \cite{loggAutomatedSolutionDifferential2012}.


We solve for the system rheology by  implementing a splitting scheme \cite{chorinNumericalSolutionNavierStokes1968,loggAutomatedSolutionDifferential2012} for the \ac{ns} equations, applied  simultaneously to the \ac{mem} and \ac{bflu} flow dynamics, and demonstrate it to be numerically stable. As for the membrane shape, we generalize the \ac{al} gauge---previously proposed to describe the steady state of membrane tubes \cite{derenyiFormationInteractionMembrane2002}---to the system dynamics, by means of a time-dependent gauge fixing. Overall, this implementation allows for handling strongly turbulent \ac{mem} and \ac{bflu} flows, as well as large \ac{mem} deformations.\vspace{0.5cm}\\

The paper is structured as follows. In \cref{sec_results} we discuss the \ac{ale} method, and present two illustrative examples: The interaction between \ac{bflu} and a \ac{rig} in \cref{sec_fl_rigid}, and that between  \ac{bflu} and \ac{ela} in \cref{sec_fl_elastic}. In \cref{sec_fl_mem} we present our approach for the interaction between \ac{bflu} and \ac{mem}. In particular, in \cref{sec_fl_mem_mem,sec_fl_mem_dom,sec_fl_mem_flu}, we discuss the splitting scheme and gauge-fixing procedure for the \ac{mem}, \ac{dom} and \ac{bflu}, respectively. Finally, \cref{sec_discussion} is devoted to the discussion of the results, and future directions.
